\documentclass{article}
\usepackage[T1]{fontenc}
\usepackage{lmodern}
\usepackage{graphicx}
\usepackage{amsmath,amssymb}
\usepackage[a4paper,margin=2cm]{geometry}
\usepackage[absolute,overlay]{textpos}
\setlength{\TPHorizModule}{1mm}
\setlength{\TPVertModule}{1mm}
\usepackage[indonesian]{babel}
\usepackage{hyperref}
\setlength{\emergencystretch}{2em}
% unit: Halaman 32

\begin{document}

% === Halaman 32 (python/native) ===

4. Kalikan vektor kolom bit (b0, b1, …, b8)T dengan sebuah matriks 

affine sebagai berikut:

5. Selanjutnya, konversi hasil perhitungan (b’0, b’1, …, b’8)T ke dalam 

heksadesimal, menjadi elemen S-box(x, y).































































































































































=





















































0

1

1

0

0

0

1

1


1

1

1

1

1

0

0

0

0

1

1

1

1

1

0

0

0

0

1

1

1

1

1

0

0

0

0

1

1

1

1

1

1

0

0

0

1

1

1

1

1

1

0

0

0

1

1

1

1

1

1

0

0

0

1

1

1

1

1

1

0

0

0

1



'

'

'

'

'

'

'

'

7

6

5

4

3

2

1

0

7

6

5

4

3

2

1

0

b

b

b

b

b

b

b

b

b

b

b

b

b

b

b

b

\end{document}
